\documentclass[letterpaper]{article}
\usepackage{aaai2027}
\usepackage[hyphens]{url}
\usepackage{graphicx}
\usepackage{booktabs}
\usepackage{longtable}
\usepackage{amsmath}
\urlstyle{rm}
\def\UrlFont{\rm}
\setcounter{secnumdepth}{0}

\title{Supplementary Material: IDT-Calibrated Latent Transport for Style-ID Artistic Transfer}
\author{Anonymous Submission}
\affiliations{}

\begin{document}
\maketitle

\section{Overview}
This supplement consolidates the paper-facing artifact paths, protocol details, and auxiliary evaluations used by the active AAAI-27 draft. It covers four items that are only summarized in the main paper:
\begin{itemize}
\item the Distinct5-WikiArt split and the selected operating-point packets;
\item the target-pooled ArtFID recompute path used for all main Distinct5 comparisons;
\item the fixed Seedream-4.5 API protocol and repaired-750 assembly rule;
\item the non-CLIP Distinct5 style probe used as an auxiliary target-style verification path.
\end{itemize}

\section{Distinct5-WikiArt Split}
The main split uses five WikiArt domains: \texttt{Early\_Renaissance}, \texttt{Impressionism}, \texttt{Minimalism}, \texttt{Rococo}, and \texttt{Ukiyo\_e}. The current local roots are:
\begin{itemize}
\item train: \texttt{Dataset/distinct5\_512/train}
\item test: \texttt{Dataset/distinct5\_512/test}
\end{itemize}
The current paper-facing classifier probe uses all 5,000 train images and all 150 held-out test images. The evaluation packet always uses the full $5\times5$ protocol over 750 generated outputs unless explicitly marked as transfer-only post-aggregation.

\section{Main-Point Ledger}
The compact artifact ledger is stored in:
\begin{itemize}
\item \texttt{aaai2027/main\_point\_artifact\_ledger.csv}
\end{itemize}
This file ties each main manuscript point to its config or protocol file, eval root, target-pooled ArtFID recompute artifact, and auxiliary metrics table.

\begin{table*}[t]
\centering
\caption{Paper-facing operating-point ledger.}
\small
\begin{tabular}{@{}llll@{}}
\toprule
Point & Role & Config or Protocol & Main Eval Root / Artifact \\
\midrule
IDT & no-op floor & \texttt{idt\_5x5 summary} & \texttt{docs/experiments/idt\_eval\_20260602/distinct5\_512/idt\_5x5} \\
LBM-K e1 & conservative anchor & \texttt{distinct5\_512\_ema\_variant\_k\_content\_adaptive\_vq\_queue\_e3.json} & \texttt{exp/distinct5\_512\_ema\_variant\_k\_content\_adaptive\_vq\_queue\_e3\_b44\_remote/full\_eval/epoch\_0001} \\
LBM-Knee e13 & main closed point & \texttt{inmortal\_xpred\_kmanifold\_pattn\_anisostokes\_queue\_from\_pattn\_seed42\_b8a2.json} & \texttt{aaai2027/local\_eval/lbm\_knee\_e13\_artfid} \\
LBM-PS-v2 e13 & style ceiling & \texttt{inmortal\_xpred\_kmanifold\_pattn\_stokes002\_from\_pattn\_seed42\_b16.json} & \texttt{aaai2027/local\_eval/pattn\_stokes002\_e13} \\
SaMAM-2250 & negative compact baseline & \texttt{2026-06-04-samam-distinct5-valid-through-2250.md} & \texttt{artfid\_comparison\_points.remote bundle + visual alignment artifact} \\
SaMST e15 & high-style baseline & \texttt{samst\_distinct5\_512\_real\_b2\_e15\_20260602/epoch\_0015 summary} & \texttt{Related\_Works/baseline\_pipeline/results/.../epoch\_0015} \\
Seedream-4.5 & external large-prior reference & \texttt{assembly\_manifest.json} & \texttt{seedream45\_api/distinct5\_512\_seedream45\_windhub\_20260607\_repaired750} \\
\bottomrule
\end{tabular}
\end{table*}

\section{Operating-Point Mechanism Table}
The main paper now defines the paper-facing variants in words. Table~\ref{tab:variant_mech} gives the compact mechanism version.

\begin{table}[t]
\centering
\caption{Compact mechanism ladder for the current paper-facing LBM rows.}
\label{tab:variant_mech}
\small
\begin{tabular}{@{}l l l l@{}}
\toprule
Variant & Endpoint selection & Texture/statistics & Motion control \\
\midrule
LBM-K & queue-guided OMF & SA-SWD & kinetic only \\
LBM-Knee & barycentric top-8 + lowfreq/edge & cross-attn proximal & anisotropic + Stokes + queue smoothing \\
LBM-PS & barycentric top-8 + lowfreq/edge & cross-attn proximal & Stokes 0.05 \\
LBM-PS-v2 & barycentric top-8 + lowfreq/edge & cross-attn proximal & Stokes 0.02 \\
\bottomrule
\end{tabular}
\end{table}

\section{Target-Pooled ArtFID Recompute}
The main paper no longer mixes the older pair-mean JSON view with the current target-pooled reading. The current fast recompute script is:
\begin{itemize}
\item \texttt{tools/compute\_targetwise\_artfid\_fast.py}
\end{itemize}
Its contract is:
\begin{enumerate}
\item pool generated images by target style;
\item recompute the art-domain FID term on the pooled target bucket;
\item recompute the LPIPS-Alex content term on the same pooled scope;
\item form per-target ArtFID as $(1+\mathrm{FID}_{art})(1+\mathrm{LPIPS}_{content})$;
\item report the mean of the five per-target ArtFID values.
\end{enumerate}

This reproduces the trusted Distinct5 reference packets to rounding tolerance:
\begin{itemize}
\item IDT all-pairs: \texttt{216.5}
\item SaMST e15 all-pairs: \texttt{395.7}
\end{itemize}

The current paper-facing recompute artifacts are:
\begin{itemize}
\item \texttt{aaai2027/local\_eval/lbm\_k\_e1\_artfid\_fast\_repro.json}
\item \texttt{aaai2027/local\_eval/lbm\_knee\_e13\_artfid/aggregate\_targetwise\_artfid\_fast\_repro.json}
\item \texttt{aaai2027/local\_eval/lbm\_psv2\_e13\_artfid\_fast\_repro.json}
\item \texttt{aaai2027/tmp\_samst\_artfid\_fast\_repro3.json}
\item \texttt{aaai2027/local\_eval/seedream\_repaired750\_artfid/aggregate\_targetwise\_artfid\_fast\_repro.json}
\end{itemize}

\section{Fixed-Rule Follow-Up Splits}
Two reviewed fixed-rule follow-up WikiArt stress splits are already closed and safe to cite as narrow generalization support:
\begin{itemize}
\item \texttt{docs/experiments/2026-06-06-faraday-split1-paper-safe-packet.md}
\item \texttt{docs/experiments/2026-06-06-faraday-split2-paper-safe-packet.md}
\end{itemize}

Closed transfer-only readout:
\begin{itemize}
\item split1 IDT: \texttt{0.6858}; split1 LBM-F e1: \texttt{0.6997}; gain over IDT: \texttt{+0.0139}
\item split2 IDT: \texttt{0.7168}; split2 LBM-F e1: \texttt{0.7241}; gain over IDT: \texttt{+0.0073}
\end{itemize}

These packets are intentionally narrow:
\begin{itemize}
\item they show that the unchanged-image pathology persists beyond Distinct5-WikiArt;
\item they show that a compact retained LBM point still clears the split-local IDT floor;
\item they do not support broader baseline or artifact-superiority claims on those follow-up splits.
\end{itemize}

\section{Seedream-4.5 Protocol}
Seedream-4.5 is treated strictly as an external large-prior reference, not as a same-interface style-ID baseline. The fixed evaluated package is:
\begin{itemize}
\item \texttt{Related\_Works/baseline\_pipeline/results/seedream45\_api/distinct5\_512\_seedream45\_windhub\_20260607\_repaired750}
\end{itemize}

Protocol summary:
\begin{itemize}
\item provider: \texttt{windhub.cc}
\item model id: \texttt{doubao-seedream-4-5-251128}
\item inference interface: source image plus a fixed target-domain prompt template
\item repaired assembly: the first full pass yielded 720 images; 30 blocked jobs were repaired using same-style replacement sources
\item replacement map: \texttt{docs/experiments/2026-06-07-distinct5-seedream45-replacements.json}
\item repaired source root for pooled content comparisons: \texttt{tmp/distinct5\_seedream\_repaired\_eval\_20260607}
\end{itemize}

The current same-scope target-pooled recompute yields:
\begin{itemize}
\item all-pairs ArtFID: \texttt{311.5}
\item transfer-only ArtFID: \texttt{361.6}
\end{itemize}

\section{Non-CLIP Style Probe}
The non-CLIP style verification path uses the existing ConvNeXt-Tiny image classifier trainer:
\begin{itemize}
\item script: \texttt{src/utils/classify.py}
\item output checkpoint: \texttt{aaai2027/distinct5\_convnext\_style\_classifier.pt}
\item output report: \texttt{aaai2027/distinct5\_convnext\_style\_classifier\_report.json}
\end{itemize}

Training setup:
\begin{itemize}
\item train root: \texttt{Dataset/distinct5\_512/train}
\item val root: \texttt{Dataset/distinct5\_512/test}
\item image size: \texttt{224}
\item batch size: \texttt{40}
\item backbone: pretrained ConvNeXt-Tiny
\end{itemize}

Held-out test result:
\begin{itemize}
\item accuracy: \texttt{0.9600}
\item macro recall: \texttt{0.9600}
\item mean confidence: \texttt{0.9643}
\end{itemize}

The paper-facing operating-point scoring script is:
\begin{itemize}
\item \texttt{tools/eval\_nonclip\_style\_probe.py}
\end{itemize}

Outputs:
\begin{itemize}
\item \texttt{aaai2027/distinct5\_nonclip\_style\_probe.csv}
\item \texttt{aaai2027/distinct5\_nonclip\_style\_probe.json}
\end{itemize}

\begin{table}[t]
\centering
\caption{Transfer-only non-CLIP target-style verification.}
\small
\begin{tabular}{@{}lrrrr@{}}
\toprule
Point & Target Acc. & Target Prob. & Source Prob. & Margin \\
\midrule
IDT & 0.010 & 0.0168 & 0.9329 & -0.9161 \\
LBM-K & 0.167 & 0.1524 & 0.7069 & -0.5544 \\
LBM-Knee & 0.237 & 0.2123 & 0.5633 & -0.3511 \\
LBM-PS-v2 & 0.272 & 0.2696 & 0.3064 & -0.0368 \\
SaMST e15 & 0.248 & 0.2405 & 0.5953 & -0.3548 \\
Seedream-4.5 & 0.378 & 0.3758 & 0.4774 & -0.1016 \\
\bottomrule
\end{tabular}
\end{table}

\section{Row-Resampled IDT Bootstrap}
The current paper keeps row-resampled transfer-only bootstrap as a narrow stability check rather than a full clustered uncertainty analysis. The extended artifact package is:
\begin{itemize}
\item \texttt{aaai2027/distinct5\_idt\_bootstrap\_extended.csv}
\item \texttt{aaai2027/distinct5\_idt\_bootstrap\_extended.md}
\end{itemize}

For the promoted rows, all resampled draws remain above the IDT floor:
\begin{itemize}
\item LBM-Knee e13: $\Delta_{\mathrm{idt,tr}}=0.0702$, interval $[0.0635, 0.0769]$
\item LBM-PS-v2 e13: $\Delta_{\mathrm{idt,tr}}=0.0901$, interval $[0.0826, 0.0976]$
\item Seedream-4.5: $\Delta_{\mathrm{idt,tr}}=0.0538$, interval $[0.0484, 0.0592]$
\end{itemize}
These are row-resampled intervals over paired transfer rows, not clustered source/style confidence intervals.

\section{Randomness and Seeds}
The current paper-facing Distinct5 packets use fixed seeds where the reviewed local configuration exposes them directly:
\begin{itemize}
\item the current \texttt{aaai2027} successor configs use seed 42 and encode this in the config names;
\item the non-CLIP ConvNeXt style probe uses seed 42;
\item the compact LBM-K anchor is referenced through its stored resolved config and eval root under the original run directory;
\item Seedream-4.5 is treated as an external API reference and does not expose a user-controlled sampling seed in the current provider interface.
\end{itemize}

\section{Reproduction Scripts}
The main paper-facing scripts referenced by the current draft are:
\begin{itemize}
\item \texttt{tools/compute\_targetwise\_artfid\_fast.py}
\item \texttt{tools/eval\_selected\_style\_metrics.py}
\item \texttt{tools/eval\_content\_edge\_purity.py}
\item \texttt{tools/eval\_nonclip\_style\_probe.py}
\item \texttt{tools/compute\_distinct5\_idt\_bootstrap.py}
\item \texttt{tools/count\_paper\_point\_params.py}
\item \texttt{aaai2027/scripts\_gen\_distinct5\_page1\_summary.py}
\item \texttt{aaai2027/scripts\_gen\_distinct5\_qualitative\_main.py}
\end{itemize}

\section{Parameter Counts}
The current paper-facing Distinct5 checkpoints can be counted with:
\begin{itemize}
\item \texttt{tools/count\_paper\_point\_params.py}
\item output: \texttt{aaai2027/paper\_point\_param\_counts.csv}
\end{itemize}

Current reviewed counts:
\begin{itemize}
\item LBM-K: \texttt{1.982M}
\item LBM-Knee: \texttt{1.946M}
\item LBM-PS-v2: \texttt{1.946M}
\end{itemize}

These counts are for the current Distinct5-WikiArt packet family. They are reported separately from the preserved historical CycleGAN-256 row in the main paper.

\section{Qualitative Provenance Note}
The current main-paper qualitative strip uses two Distinct5 cases that are available across all current paper-facing families. For the SaMAM column, the source image is a crop from the existing \texttt{distinct5\_visual\_alignment\_grid.jpg} artifact because the validated SaMAM-2250 raw image packet is remote-only in the current workspace state. The quantitative manuscript boundary remains the validated SaMAM-2250 packet.

\section{Additional Qualitative Gallery}
The supplement adds six aligned Distinct5 cases that keep the current paper-facing method set:
\texttt{Source | IDT | SaMAM-2250 | SaMST | Seedream-4.5 | LBM-K | LBM-Knee | LBM-PS-v2 | target ref}.
These rows are intended to show a broader spread than the two main-paper cases and to expose both conservative and failure-like operating regimes.

\begin{figure*}[t]
\centering
\includegraphics[width=\textwidth]{figures/fig_distinct5_qualitative_appendix_a.png}
\caption{Additional aligned Distinct5 qualitative cases, rows 1--3.}
\end{figure*}

\begin{figure*}[t]
\centering
\includegraphics[width=\textwidth]{figures/fig_distinct5_qualitative_appendix_b.png}
\caption{Additional aligned Distinct5 qualitative cases, rows 4--6.}
\end{figure*}

\section{Blind Pairwise Packet}
A blinded A/B comparison packet has been prepared for the highest-value paper comparisons:
\begin{itemize}
\item LBM-Knee vs SaMST
\item LBM-Knee vs Seedream
\item LBM-PS-v2 vs SaMST
\item LBM-K vs IDT
\end{itemize}

Bundle contents:
\begin{itemize}
\item \texttt{aaai2027/blind\_pairwise\_v1/blind\_pairwise\_manifest.csv}
\item \texttt{aaai2027/blind\_pairwise\_v1/blind\_pairwise\_rubric.md}
\item \texttt{aaai2027/blind\_pairwise\_v1/exploratory\_blind\_audit.csv}
\item \texttt{aaai2027/blind\_pairwise\_v1/exploratory\_blind\_audit.md}
\item \texttt{aaai2027/blind\_pairwise\_v1/exploratory\_blind\_audit\_summary.csv}
\item \texttt{aaai2027/blind\_pairwise\_v1/panels/}
\end{itemize}

Each panel shows \texttt{Source | Candidate A | Candidate B | Target ref}, with randomized A/B order stored only in the manifest. This packet is ready for human or external VLM scoring. Local exploratory blind rubric notes and their summaries are kept only as internal supplementary support and are not promoted as a main paper claim surface.

\section{Software and Hardware}
Current local analysis environment:
\begin{itemize}
\item Python: \texttt{3.12.0}
\item GPU: \texttt{NVIDIA GeForce RTX 4070 Laptop GPU}, 8\,GiB VRAM
\item driver: \texttt{610.47}
\item package snapshot: \texttt{aaai2027/local\_analysis\_requirements.txt}
\item environment template: \texttt{aaai2027/local\_analysis\_environment.yml}
\end{itemize}
Remote training timings reported in the main paper come from the reviewed single-lane packet logs and are not recomputed by this supplement.

\end{document}
